% TOP_PROBLEM: {"id": 30007307, "title": "Jiang–Su absorption from Cuntz-semigroup regularity", "category_id": 8, "category": {"id": 8, "name": "analysis", "display_name": "Analysis", "description": "Limits, continuity, calculus, and function theory.", "slug": "analysis", "order_index": 8, "created_at": "2026-07-31T15:26:25.671Z"}, "status": "open"}
% FIRST_POSTED: 2026-10-06
% ENTRY_KIND: statement_only
% !TeX program = lualatex
\documentclass[11pt]{article}
\usepackage[margin=25mm]{geometry}
\usepackage{fontspec}
\setmainfont{Latin Modern Roman}
\usepackage{amsmath,amssymb,mathtools}
\usepackage{unicode-math}
\setmathfont{Latin Modern Math}
\usepackage{xurl,hyperref}
\hypersetup{colorlinks=true,urlcolor=blue}
\setcounter{secnumdepth}{0}
\setlength{\parindent}{0pt}
\setlength{\parskip}{5pt}
\setlength{\emergencystretch}{3em}
\newcommand{\R}{\mathbb R}
\newcommand{\N}{\mathbb N}
\newcommand{\Z}{\mathbb Z}
\newcommand{\Q}{\mathbb Q}
\newcommand{\C}{\mathbb C}
\newcommand{\D}{\mathbb D}
\newcommand{\F}{\mathbb F}
\newcommand{\sref}[2]{\hyperlink{source:#1:#2}{[S#2]}}
\newcommand{\cataloguescope}[1]{#1}
\begin{document}
\section*{Jiang–Su absorption from Cuntz-semigroup regularity}\label{30007307}
\subsection{Definitions and mathematical statement}
\paragraph{Statement recorded in the supplied source.}
Let A be a separable nuclear complex C*-algebra. Suppose that, for every $x,y\in Cu(A)$ and integer $l\geq 1$, $(l+1)x\leq ly$ implies $x\leq y$, and for every $x\in Cu(A)$ and integer $N\geq 1$ there is one $u\in Cu(A)$ satisfying $Nu\leq x\leq (N+1)u$. Must A absorb Z? These are the full ordered-semigroup conditions, allowing zero and infinite functionals and arbitrary ideals.
\subsection{Short English statement}
Let A be a separable nuclear complex C*-algebra. Suppose that, for every $x,y\in Cu(A)$ and integer $l\geq 1$, $(l+1)x\leq ly$ implies $x\leq y$, and for every $x\in Cu(A)$ and integer $N\geq 1$ there is one $u\in Cu(A)$ satisfying $Nu\leq x\leq (N+1)u$. Must A absorb Z? These are the full ordered-semigroup conditions, allowing zero and infinite functionals and arbitrary ideals.
\subsection{Sources}
\begin{itemize}
\item[S1] Original question source or mathematical context.
\url{https://arxiv.org/abs/2506.10902v2}.
\item[S2] OpenAI, manuscript citation page. Source of the formulation motivating this local problem.
\url{https://github.com/openai/math/blob/adc7f1241b42e322a6451854ab7e4b4c146bf78a/preprints/Cuntz-comparison-and-Jiang-Su-absorption-September-23-2026/README.md}.
\end{itemize}
\end{document}
